\subsection{交换幺代数的特征标}

定义 4. --- 设 A 是 K 上的一个交换幺代数。幺特征标是指从 A 到 K 的一个保单位元态射。

在不致混淆时，可简称特征标以代替幺特征标。A 的幺特征标所成的集合记作 X(A)。若 A 是零代数，则 X(A) 为空集。

设 A 与 B 是 K 上的交换幺代数，h 是从 A 到 B 的一个保单位元态射。映射 \(\chi \mapsto \chi \circ h\)（它把 \(\mathsf{X}(\mathsf{B})\) 映到 \(\mathsf{X}(\mathsf{A})\) 中）记作 \(\mathsf{X}(h)\)。若 k 是从 B 到某个交换幺代数的态射，则有 \(\mathsf{X}(k \circ h) = \mathsf{X}(h) \circ \mathsf{X}(k)\)。映射 \(\mathsf{X}(\mathrm{Id}_{\mathrm{A}})\) 是 \(\mathsf{X}(\mathrm{A})\) 上的恒同映射。

若 h 是满射，则 \(\mathsf{X}(h)\) 是从 \(\mathsf{X}(\mathsf{B})\) 到在 h 的核上为零的 A 的特征标所成集合上的一个双射。

设 \((\mathrm{A}_{1},e_{1}),\ldots,(\mathrm{A}_{n},e_{n})\) 是 K 上的交换幺代数，A 是幺代数 \(A_{1}\times\cdots\times A_{n}\)，其单位元为 \((e_{1},\ldots,e_{n})\)。对每个 i，把 \(A_{i}\) 等同于 A 的一个理想，并设 \(\pi_{i}\) 是从 A 到 \(A_{i}\) 上的典范映射。则 \(\mathsf{X}(\pi_{i})\) 是从 \(\mathsf{X}(\mathrm{A}_{i})\) 到集合 \(X_{i}\)（即 A 的在 \(\prod_{j\neq i}A_{j}\) 上为零的特征标所成的集合）上的一个双射。诸集合 \(X_{i}\) 两两不相交。另一方面，设 \(\chi\in\mathsf{X}(\mathrm{A})\)。由于 \(1=\sum\chi(e_{i})\)，存在 i 使 \(\chi(e_{i})\neq0\)。对每个 \(j\neq i\) 及每个 \(y\in A_{j}\)，有 \(\chi(e_{i})\chi(y)=\chi(e_{i}y)=\chi(0)=0\)，因而 \(\chi(\mathrm{A}_{j})=0\)。于是 \(\chi\) 在 \(\prod_{j\neq i}A_{j}\) 上为零，从而 \(\mathsf{X}(\mathrm{A})\) 是诸 \(X_{i}\) 的并。

设 B 是幺代数 \(A_{1} \otimes \cdots \otimes A_{n}\)。用 \(h_{i}\) 表示典范态射 \(A_{i} \to B\)。则

\[
\chi \mapsto (\chi \circ h _ {1}, \dots , \chi \circ h _ {n})
\]

是从 \(\mathsf{X}(\mathsf{B})\) 到 \(\mathsf{X}(\mathsf{A}_1)\times \dots \times \mathsf{X}(\mathsf{A}_n)\) 中的一个映射，而

\[
\left(\chi_ {1}, \dots , \chi_ {n}\right) \mapsto \chi_ {1} \otimes \dots \otimes \chi_ {n}
\]

是从 \(\mathsf{X}(\mathsf{A}_{1}) \times \cdots \times \mathsf{X}(\mathsf{A}_{n})\) 到 \(\mathsf{X}(\mathsf{B})\) 中的一个映射。可以验证这两个映射互为逆双射，据此我们把 \(\mathsf{X}(\mathsf{B})\) 与 \(\mathsf{X}(\mathsf{A}_{1}) \times \cdots \times \mathsf{X}(\mathsf{A}_{n})\) 视为同一。

设 A 是 K 上的一个交换幺代数。设 Y 是 A 的余维数为 1 的理想所成的集合。对每个 \(\chi \in \mathsf{X}(\mathsf{A})\)，有 \(\operatorname{Ker}(\chi) \in \mathsf{Y}\)。映射 \(\chi \mapsto \operatorname{Ker}(\chi)\) 是从 \(\mathsf{X}(\mathsf{A})\) 到 Y 上的一个双射。事实上，若 \(I \in Y\)，则存在唯一一个从幺 K-代数 A/I 到 K 上的同构，而复合态射

\[
\mathrm{A} \longrightarrow \mathrm{A} / \mathrm{I} \longrightarrow \mathrm{K}
\]

是 A 的以 I 为核的唯一特征标。

定义 5. --- 设 A 是 K 上的一个交换幺代数。对每个 \(x \in \mathrm{A}\)，我们用 \(\mathcal{G}_{\mathrm{A}}(x)\)（或简作 \(\mathcal{G}(x)\)）表示映射 \(\chi \mapsto \chi(x)\)（从 \(\mathrm{X}(\mathrm{A})\) 到 K）。它称为 \(x\) 的 Gelfand 变换。

映射 \(\mathcal{G}\) 是从 A 到幺代数 \(\mathrm{K}^{\mathsf{X}(A)}\)（由 \(\mathsf{X}(A)\) 到 K 的映射所成的幺代数）中的一个保单位元态射。它称为 A 的 Gelfand 变换。

若 \(x \in \mathrm{A}\)，则 Gelfand 变换 \(\mathcal{G}_{\mathrm{A}}(x)\)（即 \(x\) 的 Gelfand 变换）的像含于 \(\mathrm{Sp}_{\mathrm{A}}(x)\)。事实上，设 \(\chi \in \mathsf{X}(\mathrm{A})\)；由于 \(\chi(x - \chi(x)e) = 0\)，元素 \(x - \chi(x)e\) 不可逆。

设 B 是 K 上的一个交换幺代数，h 是从 A 到 B 的一个保单位元态射；于是 \(\mathsf{X}(h)\colon\mathsf{X}(\mathsf{B})\to\mathsf{X}(\mathsf{A})\) 定义了一个保单位元态射 \(h_{*}\colon\mathsf{K}^{\mathsf{X}(\mathsf{A})}\to\mathsf{K}^{\mathsf{X}(\mathsf{B})}\)，且图表：

\[
\begin{array}{l} \text {A} \xrightarrow {\mathcal {G} _ {\text {A}}} \text {K} ^ {\mathsf {X} (\text {A})} \\ \Big \downarrow h \qquad \qquad \qquad \Big \downarrow h _ {*} \\ \text {B} \xrightarrow {\mathcal {G} _ {\text {B}}} \text {K} ^ {\mathsf {X} (\text {B})} \end{array}\tag{3}
\]

是交换的。事实上，对每个 \(x \in A\) 与每个 \(\chi \in \mathsf{X}(\mathsf{B})\)，我们有：

\[
\begin{array}{r l} & {\mathcal {G} _ {\mathrm{B}} (h (x)) (\chi) = \chi (h (x)) = (\chi \circ h) (x)} \\ & {\qquad = (\mathsf {X} (h) (\chi)) (x)} \\ & {\qquad = \mathcal {G} _ {\mathrm{A}} (x) (\mathsf {X} (h) (\chi))} \\ & {\qquad = h _ {*} (\mathcal {G} _ {\mathrm{A}} (x)) (\chi).} \end{array}\tag{4}
\]

现在设 K 是一个拓扑域。这时我们给 \(\mathsf{X}(\mathsf{A})\) 赋予 A 上逐点收敛的拓扑（参见 EVT, III, p.~14, 例 1），而拓扑空间 \(\mathsf{X}(\mathsf{A})\) 称为 A 的特征标空间。于是 \(\mathsf{X}(\mathsf{A})\) 的拓扑是使诸函数 \(\mathcal{G}_{\mathrm{A}}(x)\)（\(x \in A\)）都连续的最粗拓扑，且映射 \(\chi \mapsto (\chi(a))_{a \in \mathrm{A}}\) 把空间 \(\mathsf{X}(\mathsf{A})\) 等同于 \(K^{A}\) 的一个子集。

当 K = R 或 C 时，这个拓扑是在 \(\mathsf{X}(\mathsf{A}) \subset \mathsf{A}^{*}\) 上由弱拓扑 \(\sigma(\mathsf{A}^{*}, \mathsf{A})\)（定义在 \(A^{*}\) 上）诱导的拓扑 (EVT, II, p.~45, 定义 2)；据此，我们也称它为 \(\mathsf{X}(\mathsf{A})\) 上的弱拓扑。

若 \(h\) 是从 A 到 B 的一个保单位元态射，则映射 \(\mathsf{X}(h)\colon \mathsf{X}(\mathsf{B})\to \mathsf{X}(\mathsf{A})\) 连续。若 \(h\) 是满射，则 \(\mathsf{X}(h)\) 的像，即在 \(h\) 的核上为零的 A 的特征标所成的集合，在 \(\mathsf{X}(\mathsf{A})\) 中是闭的；另一方面，\(\mathsf{X}(h)(\mathsf{X}(\mathsf{B}))\) 上由 \(\mathsf{X}(\mathsf{B})\) 的拓扑经由双射 \(\mathsf{X}(h)\) 诱导而得的拓扑，是 A 上逐点收敛的拓扑，也就是由 \(\mathsf{X}(\mathsf{A})\) 的拓扑诱导的拓扑。换言之，\(\mathsf{X}(h)\) 是从 \(\mathsf{X}(\mathsf{B})\) 到 \(\mathsf{X}(\mathsf{A})\) 的一个闭子集上的同胚。

若 \(A_{1},\ldots,A_{n}\) 是 K 上的交换幺代数，则空间 \(\mathsf{X}(\mathsf{A}_{1}\times\cdots\times\mathsf{A}_{n})\) 就这样等同于 \(\mathsf{X}(\mathsf{A}_{1}),\ldots,\mathsf{X}(\mathsf{A}_{n})\) 的拓扑和空间。类似地，\(\mathsf{X}(\mathsf{A}_{1}\otimes\cdots\otimes\mathsf{A}_{n})\) 等同于积拓扑空间 \(\mathsf{X}(\mathsf{A}_{1})\times\cdots\times\mathsf{X}(\mathsf{A}_{n})\)。

